\begin{Solbook}{1}
 \par \noindent  a) $x\ge \dfrac {3}{2}$ \hspace {2cm} b) $-\dfrac {4}{3}$  \hfill \qedEX 
\end{Solbook}
\begin{Solbook}{2}
 \par \noindent  a) $3\sqrt {3}-3$ \hspace {2cm} b) $\dfrac {1}{\sqrt {x}+4}$  \hfill \qedEX 
\end{Solbook}
\begin{Solbook}{3}
 \par \noindent Nam: $16$;\; Nữ: $19$. \hfill \qedEX 
\end{Solbook}
\begin{Solbook}{4}
 \par \noindent $20$ (áo/giờ). \hfill \qedEX 
\end{Solbook}
\begin{Solbook}{5}
 \par \noindent $x\le -\dfrac {6}{5}$. \hfill \qedEX 
\end{Solbook}
\begin{Solbook}{6}
 \par \noindent \immini {\begin {enumerate} \item $\widehat {BOA}=60^\circ $. \\Chứng minh $\triangle OAB= \triangle OAC$ (c.c.c) \par Suy ra $\widehat {ACO}=90^\circ $ \par Do đó $AC$ là tiếp tuyến của $(O;R)$. \item Chứng minh $\triangle CDB \sim \triangle HBA$ (g.g) \par Suy ra $\widehat {CDB}=\widehat {ABH}$. \item Dùng định lý Pythagore trong các tam giác vuông $MBH$ và $MHE$, chứng minh $MB^{2}-ME^{2}=HB^2-HE^2$. Sau đó dùng $HB=HC$ để biến đổi $HB^2-HE^2=EB\cdot EC$ \end {enumerate}} {\begin {tikzpicture}[scale=1, font=\footnotesize ] \coordinate (O) at (0,0); \def \R {1.7} \draw (O) circle(\R ); \coordinate (A) at (-2*\R ,0); \path [name path=cirA] (-\R ,0) circle(\R ); \path [name path=cirO] (O) circle(\R ); \path [name intersections={of=cirO and cirA, by={B,C}}]; \coordinate (D) at ($(B)!2!(O)$); \coordinate (H) at ($(B)!1/2!(C)$); \coordinate (M) at (180:\R ); \coordinate (E) at (intersection of M--D and B--C); \draw (A)--(B)--(C)--(A)--(M)--(B)--(D)--(C)--(O)--(M)--(D); \foreach \x /\g in{O/0,A/-180,B/90,C/-135,H/45,D/-90,M/135,E/-135} {\draw [fill=white] (\x ) circle (1.5pt) +(\g :.3) node{$\x $};} \end {tikzpicture}}  \hfill \qedEX 
\end{Solbook}
\begin{Solbook}{7}
 \par \noindent Diện tích phần tô màu là $18{,}27\,\text {cm}^2$. \hfill \qedEX 
\end{Solbook}
